% =============================================================
% BODY (EN) – GENERAL TEMPLATE CONTENT
% =============================================================

\section{Overview}
\vspace{2mm}

This document uses a bilingual (ES/EN) template for technical and personal documentation. The goal is to keep the structure modular, consistent, and easy to maintain over time.

\vspace{2mm}


\subsection{What To Edit}
\vspace{2mm}

\begin{itemize}
  \item \textbf{Metadata:} edit the \texttt{METADATA} block (title, category, version, status, date, author).
  \item \textbf{Content:} write in \texttt{lang/en/body.tex}.
  \item \textbf{References:} update \texttt{src/references.tex} (manual bibliography at the end).
  \item \textbf{Assets:} store resources in \texttt{assets/figures/}, \texttt{assets/diagrams/}, and \texttt{assets/tables/}.
\end{itemize}

\vspace{2mm}


\section{Main Content}
\vspace{2mm}

Keep paragraphs short and introduce subsections as the topic grows (one main idea per subsection).

\vspace{2mm}


\subsection{Notes, Tips, And Warnings}
\vspace{2mm}

\begin{noteBox}
\boxlabel{Note}
Use this block for context, decisions, or important clarifications.
\end{noteBox}

\vspace{2mm}

\begin{tipBox}
\boxlabel{Tip}
Use this block for practical recommendations and safe shortcuts.
\end{tipBox}

\vspace{2mm}

\begin{warnBox}
\boxlabel{Warning}
Use this block for risks, constraints, or error-prone steps.
\end{warnBox}

\vspace{2mm}


\subsection{Checklist}
\vspace{2mm}

\begin{itemize}
  \item Item to verify.
  \item Assumption or decision taken.
  \item Expected outcome.
  \item Next step.
\end{itemize}

\vspace{2mm}


\section{Examples}
\vspace{2mm}

\subsection{Code}
\vspace{2mm}

\begin{lstlisting}
# Example
print("Hello, world!")
\end{lstlisting}

\vspace{2mm}


\subsection{Figure (Snippet)}
\vspace{2mm}

\begin{lstlisting}
\begin{figure}[htbp]
  \centering
  \includegraphics[width=0.80\linewidth]{assets/figures/your-figure.png}
  \caption{Short figure caption.}
\end{figure}
\end{lstlisting}

\vspace{2mm}


\subsection{Diagrams / Logic Gates (TikZ)}
\vspace{2mm}

\begin{center}
\begin{tikzpicture}[circuit logic US, scale=1.0, transform shape]
  \node[and gate, draw, logic gate inputs=nn] (G) at (0,0) {};
  \draw (-1.2,0.2) -- (G.input 1);
  \draw (-1.2,-0.2) -- (G.input 2);
  \draw (G.output) -- (1.2,0);
\end{tikzpicture}
\end{center}

\vspace{2mm}